% Generated draft: requires mathematical and source review.
\documentclass[11pt]{article}
\usepackage[margin=25mm]{geometry}
\usepackage{fontspec}
\setmainfont{Latin Modern Roman}
\usepackage{amsmath,amssymb,mathtools,unicode-math}
\setmathfont{Latin Modern Math}
\usepackage{enumitem,needspace,xurl,hyperref}
\setlength{\parindent}{0pt}
\setlength{\parskip}{5pt}
\setlength{\emergencystretch}{3em}
\newcommand{\N}{\mathbb{N}}
\newcommand{\Z}{\mathbb{Z}}
\newcommand{\Q}{\mathbb{Q}}
\newcommand{\R}{\mathbb{R}}
\newcommand{\C}{\mathbb{C}}
\newcommand{\F}{\mathbb{F}}
\newcommand{\A}{\mathbb{A}}
\newcommand{\PP}{\mathbb{P}}
\newcommand{\E}{\mathbb{E}}
\newcommand{\eps}{\varepsilon}
\newcommand{\dd}{\,\mathrm{d}}
\begin{document}
\section*{UnsolvedMath ID: 30002281}
\textbf{Source problem code:} catalogue entry
\subsection{Definitions and mathematical statement}
The supplied record calls this the motivic stabilization conjecture and assumes that \(X\) is geometrically irreducible. It does not specify the base field, the category of varieties, or the topological ring in which convergence is asserted. These omissions prevent a fully specified convergence statement.

In the customary motivic notation, for a field \(k\), geometric irreducibility means that \(X\times_k\overline{k}\) is irreducible. The symmetric power is \(\operatorname{Sym}^nX=X^n/\mathfrak S_n\), when this quotient exists in the chosen category, where \(\mathfrak S_n\) permutes the factors. The Grothendieck ring \(K_0(\mathrm{Var}_k)\) is generated by classes \([Y]\) of \(k\)-varieties, with relations \([Y]=[Z]+[Y\setminus Z]\) for closed subvarieties \(Z\subseteq Y\), and multiplication induced by products over \(k\). Write \(\mathbb L=[\mathbb A^1_k]\). For \(d=\dim X\), the normalized classes are
\[
 a_n(X)=[\operatorname{Sym}^nX]\mathbb L^{-nd}
 \in K_0(\mathrm{Var}_k)[\mathbb L^{-1}].
\]
These definitions explain the customary notation; the record does not explicitly supply this ambient setup.

The source assertion is that, for every geometrically irreducible \(X\) in the intended domain,
\[
 \lim_{n\to\infty}\frac{[\operatorname{Sym}^nX]}{\mathbb L^{n\dim X}}
 \quad\text{exists}.
\]
More precisely, once a topological ring \(R\) and the intended homomorphism \(\iota:K_0(\mathrm{Var}_k)[\mathbb L^{-1}]\to R\) are specified, this means that for every such \(X\) there exists \(\ell_X\in R\) such that, for every neighborhood \(U\) of \(0\) in \(R\), there is an integer \(N\geq1\) for which
\[
 n\geq N\ \Longrightarrow\ \iota(a_n(X))-\ell_X\in U.
\]
The intended ring, topology, any completion or quotient, and the allowed fields and varieties remain unresolved; this parametrized explanation does not choose them on behalf of the source.
\subsection{Short English statement}
The proposed assertion is that the motivic classes of the symmetric powers of every geometrically irreducible \(X\), normalized by \(\mathbb L^{n\dim X}\), converge as \(n\to\infty\). The supplied formulation leaves the base field and the ambient topology unspecified, so the exact convergence claim requires clarification. Its open status is supplied dated curation, not an independently verified conclusion.
\subsection{Sources}
\begin{itemize}
\item Local UnsolvedMath registry, permanent ID \(30002281\), title \emph{Motivic Stabilization of Symmetric Powers}, source problem code \texttt{catalogue entry}, in the Oberwolfach Reports Open Problems collection. The record gives the proposed year as 2013 and supplies no named proposer.
\item Supplied report citation: \emph{Moduli Spaces in Algebraic Geometry} (2013), Oberwolfach Reports, pp.~343--392, DOI \texttt{10.4171/owr/2013/06}. \url{https://doi.org/10.4171/owr/2013/06}. The supplied background describes Oberwolfach Reports as documenting mathematical workshops and their open-problem discussions.
\item Supplied literature reference: Vakil--Wood, \emph{discriminants in the Grothendieck ring}. \url{https://arxiv.org/abs/1208.3166}.
\item Supplied literature reference: Litt, \emph{symmetric powers do not stabilize}. \url{https://arxiv.org/abs/1209.4708}.
\item Literature assessment dated exactly \texttt{2026-08-21}: the supplied curation labels the problem open and states that the Vakil--Wood motivic stabilization conjecture remains unresolved in general, with related nonstabilization results not settling this exact normalized limit. This is supplied dated literature triage derived from the OWR curation, not independently verified here; its time-sensitive status and cited claims require independent verification.
\item The registry labels the statement \texttt{corrected\_verified} and reports that its mathematical content was checked against the official Oberwolfach report, with notation, context, and question form normalized without changing the source-level claim. This verification claim is supplied metadata, not an independent verification performed for this draft.
\end{itemize}
\end{document}
